\chapter{Target-Population Sensitivity and Regression Results}
\label{ch:modeling_strategy_results}

\section{Empirical Estimand Ladder}

This chapter is the empirical core of the thesis. It applies the framework defined in Chapter~\ref{ch:framework} to the Eurostat unmet-care panels. The poverty or social-exclusion coefficient is treated as a stress-test coefficient: it is followed as the model moves from the intended outcome-observed monitoring panel to selected complete cases, population-weighted complete cases, universe-restricted panels, IPW diagnostics, MAR imputation, and MNAR stress tests. The question is not whether one selected-panel coefficient is statistically marked. The question is whether the aggregate poverty association is stable when the target population, weighting rule, row-selection process, and missing-data assumption change.

\begin{table}[htbp]
\centering
\scriptsize
\caption{Target populations, assumptions, and interpretation of estimates. Row counts come from the generated thesis pipeline and distinguish the 608-row outcome-observed target from the 282-row complete-case FE sample.}
\label{tab:estimand_framework}
\resizebox{\textwidth}{!}{%
\begin{tabular}{@{}p{0.16\textwidth}p{0.25\textwidth}rrp{0.20\textwidth}p{0.25\textwidth}@{}}
\toprule
Label & Target population & Countries & Rows & Weighting & Interpretation \\
\midrule
Outcome-observed & All country-years with observed \texttt{hlth\_silc\_08} & 37 & 608 & None & Average-country monitoring target before covariate attrition. \\
Population-weighted & Outcome-observed or selected rows with valid \texttt{demo\_pjan} & 37 & 607 & Year-normalized population & Average-resident monitoring target; not inherently more correct. \\
Complete-case FE & Outcome and all baseline regression covariates observed & 30 & 282 & None & Selected analytical panel for the baseline FE coefficient. \\
Balanced-outcome & Countries with complete primary-outcome and population-weight coverage & 25 & 238 & Optional & Coverage-restricted target; not automatically covariate-balanced. \\
IPW CC & Complete cases reweighted by predicted complete-case inclusion & 30 & 282 & Stabilized truncated IPW & Selection diagnostic; cannot recover structurally missing covariates. \\
MAR MI & Outcome-observed rows with model-imputed covariates & 37 & 608 & Usually none & Model-based sensitivity target under imputation assumptions. \\
MNAR sensitivity & Outcome-observed rows with shifted imputed values & 37 & 608 & Scenario-specific & Stress test for departures from MAR; implemented in Step 2. \\
\bottomrule
\end{tabular}%
}

\vspace{0.2em}
\begin{minipage}{0.96\textwidth}\scriptsize Notes: Row counts refer to the generated thesis pipeline. Balanced-outcome means balanced for the primary outcome and population weights, not for every regression covariate. MAR and MNAR imputation estimates are model-based sensitivity estimates.\end{minipage}

\vspace{0.25em}

\begin{minipage}{0.96\textwidth}
\scriptsize \textit{Source:} Author's synthesis from the thesis data pipeline and missing-data literature \citep{little_rubin_2019,seaman_white_2013}. The table is a conceptual map for interpreting generated estimates, not an external empirical result.
\end{minipage}
\end{table}

The hierarchy is estimand-first. Complete-case and population-weighted complete-case estimates use the same selected rows but answer different country-versus-resident monitoring questions. Balanced-outcome means balanced for the primary outcome and valid population weights, not covariate-balanced. IPW is a selection diagnostic among observed complete cases. MAR and MNAR MI rows target the broader outcome-observed panel under explicit model-based assumptions and are not treated as corrected estimates \citep{rubin_1987,little_rubin_2019}.

\section{Multi-Outcome Target-Population Drift Diagnostics}

Before estimating poverty/social-exclusion coefficients, the audit protocol asks how far each feasible outcome can travel from the monitoring target to an analytical sample. Table~\ref{tab:multi_outcome_attrition_matrix} reports this drift across the multi-outcome registry. The columns are sequential requirements: outcome observed, valid population weight, GDP, unemployment, poverty/social exclusion, the baseline financing-complete covariate set, and a stricter capacity-augmented sensitivity set.

The matrix shows that target-population drift is not a one-off feature of the primary medical indicator. Long population-denominator medical and dental indicators have broad monitoring coverage, but their baseline complete-case analytical samples are much smaller. Need-denominator indicators are shorter by construction and become especially limited once financing and capacity covariates are required. These rows do not estimate any substantive effect; they define which empirical targets are supportable.

Figure~\ref{fig:multi_outcome_retention_heatmap} visualizes the same row-retention problem. The heatmap is useful because it separates denominator limits from covariate-completeness limits. A short need-denominator indicator may be feasible for monitoring but weak for regression once country-year covariates are required.

The detailed 608-to-282 primary-outcome sequence is reported once in Figure~\ref{fig:complete_case_attrition_waterfall} and Table~\ref{tab:complete_case_attrition_waterfall}. The multi-outcome matrix and retention heatmap extend that diagnostic across outcomes without repeating the primary waterfall.

Table~\ref{tab:multi_outcome_included_excluded_smd_top} shows that attrition is not only about row counts. Excluded rows differ from included rows on variables tied to vulnerability and system capacity. The direction and size of these differences vary by indicator length and denominator, which is why the later coefficient matrix labels infeasible or unstable estimands instead of forcing them into one preferred model.

\begin{table}[H]
\centering
\scriptsize
\caption{Multi-outcome attrition matrix. Rows are feasible Eurostat unmet-care indicators; columns report retained country-year rows after sequentially requiring population weights and covariate blocks. Baseline CC is the financing-complete fixed-effects covariate set used for the primary poverty stress test, while capacity is a stricter sensitivity requirement.}
\label{tab:multi_outcome_attrition_matrix}
\begin{tabular}{p{0.24\linewidth}rrrrrrr}
\toprule
Indicator & Outcome & Pop. & GDP & Unemp. & Poverty & Baseline CC & Capacity \\
\midrule
Dental need & 154 & 154 & 152 & 147 & 147 & 93 & 69 \\
Dental population & 607 & 606 & 604 & 555 & 360 & 282 & 233 \\
Medical need & 154 & 154 & 152 & 147 & 147 & 93 & 69 \\
Medical population & 608 & 607 & 604 & 555 & 360 & 282 & 233 \\
Medical cost & 608 & 607 & 604 & 555 & 360 & 282 & 233 \\
Medical distance & 608 & 607 & 604 & 555 & 360 & 282 & 233 \\
Medical waiting & 608 & 607 & 604 & 555 & 360 & 282 & 233 \\
\bottomrule
\end{tabular}

\vspace{0.25em}

\begin{minipage}{0.96\textwidth}
\scriptsize \textit{Source:} Author's calculations from the multi-outcome Eurostat panel and \texttt{demo\_pjan}. The table is a target-support diagnostic: it reports which outcome-by-covariate targets can support later estimands.
\end{minipage}
\end{table}

\begin{figure}[H]
\centering
\includegraphics[width=0.95\textwidth]{multi_outcome_retention_heatmap}
\caption{Multi-outcome retention heatmap. Cells report the percentage of each indicator's outcome-observed country-years retained after sequential population-weight and covariate requirements; no regression is estimated in this diagnostic.}
\label{fig:multi_outcome_retention_heatmap}
\vspace{0.25em}

\begin{minipage}{0.94\textwidth}
\scriptsize \textit{Source:} Author's calculations from the multi-outcome Eurostat panel. Retention is measured relative to each indicator's own outcome-observed country-year target.
\end{minipage}
\end{figure}

\begin{table}[H]
\centering
\scriptsize
\caption{Largest included-versus-excluded standardized differences by outcome. Included rows are baseline complete-case rows for each outcome; excluded rows are outcome-observed rows removed by missing baseline covariates. SMDs are descriptive selection diagnostics.}
\label{tab:multi_outcome_included_excluded_smd_top}
\begin{tabular}{p{0.20\linewidth}p{0.21\linewidth}rrrrr}
\toprule
Indicator & Variable & In N & Excl. N & In mean & Excl. mean & SMD \\
\midrule
Dental need & physicians\_per\_100k & 71 & 7 & 413.741 & 326.734 & 1.187 \\
Dental need & poverty\_or\_social\_exclusion\_pc & 93 & 61 & 20.304 & 24.361 & -0.63 \\
Dental population & physicians\_per\_100k & 236 & 195 & 385.83 & 325.958 & 0.907 \\
Dental population & poverty\_or\_social\_exclusion\_pc & 282 & 93 & 21.0 & 28.39 & -0.903 \\
Medical need & physicians\_per\_100k & 71 & 7 & 413.741 & 326.734 & 1.187 \\
Medical need & poverty\_or\_social\_exclusion\_pc & 93 & 61 & 20.304 & 24.361 & -0.63 \\
Medical population & physicians\_per\_100k & 236 & 195 & 385.83 & 325.958 & 0.907 \\
Medical population & poverty\_or\_social\_exclusion\_pc & 282 & 93 & 21.0 & 28.39 & -0.903 \\
Medical cost & physicians\_per\_100k & 236 & 195 & 385.83 & 325.958 & 0.907 \\
Medical cost & poverty\_or\_social\_exclusion\_pc & 282 & 93 & 21.0 & 28.39 & -0.903 \\
Medical distance & physicians\_per\_100k & 236 & 195 & 385.83 & 325.958 & 0.907 \\
Medical distance & poverty\_or\_social\_exclusion\_pc & 282 & 93 & 21.0 & 28.39 & -0.903 \\
Medical waiting & physicians\_per\_100k & 236 & 195 & 385.83 & 325.958 & 0.907 \\
Medical waiting & poverty\_or\_social\_exclusion\_pc & 282 & 93 & 21.0 & 28.39 & -0.903 \\
\bottomrule
\end{tabular}

\vspace{0.25em}

\begin{minipage}{0.96\textwidth}
\scriptsize \textit{Source:} Author's calculations from feasible Eurostat unmet-care indicators and baseline covariates. SMDs diagnose selection distortion; they are not regression coefficients.
\end{minipage}
\end{table}

\FloatBarrier

\section{Outcome-by-Estimand Coefficient Matrix}

Table~\ref{tab:outcome_estimand_coefficient_matrix} replaces a long list of separate robustness rows with a single outcome-by-estimand audit. For each feasible outcome, the table follows the poverty/social-exclusion coefficient across five estimand families: unweighted complete-case FE, population-weighted complete-case FE, balanced-outcome FE, IPW complete-case FE, and MAR MI full-target FE. Rows that do not meet the minimum support rule are reported as infeasible rather than forced.

The matrix changes the interpretation of the empirical evidence. Long population-denominator indicators can support several observed-data estimands. Need-denominator indicators are useful for monitoring, but their observed complete-case FE rows are too short for the deterministic support rule and are therefore non-identifiable as complete-case regression targets. MAR MI can still be estimated for those shorter targets, but one imputed estimand family is not enough to support a stable coefficient interpretation.

Figure~\ref{fig:outcome_estimand_coefficient_stability} shows the core empirical contribution of the chapter. Observed complete-case estimands often remain positive, while MAR MI full-target estimates can move toward zero or change sign. The result is a coefficient-stability audit: the same poverty/social-exclusion stress-test variable has different interpretations under different target populations and missing-data assumptions.

Table~\ref{tab:drift_stability_summary} translates the coefficient matrix into the framework metrics defined in Chapter~\ref{ch:framework}. The poverty/social-exclusion coefficient remains a stress-test coefficient, not the thesis contribution. The standardized coefficient \(\theta\) puts coefficients on a comparable IQR scale across outcomes; \(\theta\) IQR summarizes dispersion across feasible estimand families; directional agreement records whether signs agree with the observed-data reference direction; and CSI applies the declared \(\delta=0.10\) tolerance rule. The final labels are deterministic diagnostic labels, not causal claims and not evidence of individual-level mechanisms.

\begin{table}[H]
\centering
\scriptsize
\caption{Outcome-by-estimand coefficient matrix for the poverty/social-exclusion coefficient. Each row defines a distinct target population and missing-data assumption; infeasible rows are retained with the reason. MAR MI is interpreted as model-based sensitivity, not as correction.}
\label{tab:outcome_estimand_coefficient_matrix}
\begingroup
\scriptsize
\setlength{\tabcolsep}{2pt}
\renewcommand{\arraystretch}{0.92}
\begin{longtable}{@{}p{0.13\linewidth}p{0.16\linewidth}rrrp{0.055\linewidth}p{0.135\linewidth}p{0.055\linewidth}p{0.15\linewidth}@{}}
\caption{Outcome-by-estimand coefficient matrix for the poverty/social-exclusion coefficient. Each row defines a distinct target population and missing-data assumption; infeasible rows are retained with the reason. MAR MI is interpreted as model-based sensitivity, not as correction.}\label{tab:outcome_estimand_coefficient_matrix}\\
\toprule
Indicator & Estimand & Rows & Ctries & Years & Coef. & 95\% CI & p & Status \\
\midrule
\endfirsthead
\caption[]{Outcome-by-estimand coefficient matrix for the poverty/social-exclusion coefficient. Each row defines a distinct target population and missing-data assumption; infeasible rows are retained with the reason. MAR MI is interpreted as model-based sensitivity, not as correction. (continued)}\\
\toprule
Indicator & Estimand & Rows & Ctries & Years & Coef. & 95\% CI & p & Status \\
\midrule
\endhead
\midrule
\multicolumn{9}{r}{\scriptsize Continued on next page}\\
\endfoot
\bottomrule
\endlastfoot
Dental need & Unweighted CC FE & 93 & 28 & 4 &  &  &  & infeas.: rows<100 \\
Dental need & Pop.-weighted CC FE & 93 & 28 & 4 &  &  &  & infeas.: rows<100 \\
Dental need & Balanced FE & 86 & 24 & 4 &  &  &  & infeas.: rows<100 \\
Dental need & IPW CC FE & 93 & 28 & 4 &  &  &  & infeas.: rows<100 \\
Dental need & MAR MI full target & 154 & 34 & 5 & -0.232 & [-0.929, 0.465] & 0.514 & est. \\
Dental population & Unweighted CC FE & 282 & 30 & 10 & 0.203 & [-0.066, 0.471] & 0.139 & est. \\
Dental population & Pop.-weighted CC FE & 282 & 30 & 10 & 0.491 & [0.029, 0.953] & 0.037 & est. \\
Dental population & Balanced FE & 238 & 25 & 10 & 0.226 & [0.068, 0.383] & 0.005 & est. \\
Dental population & IPW CC FE & 282 & 30 & 10 & 0.220 & [-0.103, 0.544] & 0.182 & est. \\
Dental population & MAR MI full target & 607 & 37 & 18 & -0.009 & [-0.096, 0.079] & 0.848 & est. \\
Medical need & Unweighted CC FE & 93 & 28 & 4 &  &  &  & infeas.: rows<100 \\
Medical need & Pop.-weighted CC FE & 93 & 28 & 4 &  &  &  & infeas.: rows<100 \\
Medical need & Balanced FE & 86 & 24 & 4 &  &  &  & infeas.: rows<100 \\
Medical need & IPW CC FE & 93 & 28 & 4 &  &  &  & infeas.: rows<100 \\
Medical need & MAR MI full target & 154 & 34 & 5 & 0.094 & [-0.132, 0.320] & 0.414 & est. \\
Medical population & Unweighted CC FE & 282 & 30 & 10 & 0.178 & [0.046, 0.309] & 0.008 & est. \\
Medical population & Pop.-weighted CC FE & 282 & 30 & 10 & 0.304 & [-0.007, 0.615] & 0.055 & est. \\
Medical population & Balanced FE & 238 & 25 & 10 & 0.154 & [0.060, 0.248] & 0.001 & est. \\
Medical population & IPW CC FE & 282 & 30 & 10 & 0.214 & [0.060, 0.367] & 0.006 & est. \\
Medical population & MAR MI full target & 608 & 37 & 18 & -0.012 & [-0.098, 0.074] & 0.789 & est. \\
Medical cost & Unweighted CC FE & 282 & 30 & 10 & 0.183 & [0.088, 0.277] & <.001 & est. \\
Medical cost & Pop.-weighted CC FE & 282 & 30 & 10 & 0.299 & [0.066, 0.533] & 0.012 & est. \\
Medical cost & Balanced FE & 238 & 25 & 10 & 0.146 & [0.083, 0.208] & <.001 & est. \\
Medical cost & IPW CC FE & 282 & 30 & 10 & 0.230 & [0.099, 0.362] & <.001 & est. \\
Medical cost & MAR MI full target & 608 & 37 & 18 & 0.019 & [-0.061, 0.100] & 0.640 & est. \\
Medical distance & Unweighted CC FE & 282 & 30 & 10 & 0.008 & [-0.002, 0.018] & 0.097 & est. \\
Medical distance & Pop.-weighted CC FE & 282 & 30 & 10 & 0.018 & [0.000, 0.036] & 0.048 & est. \\
Medical distance & Balanced FE & 238 & 25 & 10 & 0.011 & [0.003, 0.020] & 0.007 & est. \\
Medical distance & IPW CC FE & 282 & 30 & 10 & 0.007 & [-0.002, 0.016] & 0.148 & est. \\
Medical distance & MAR MI full target & 608 & 37 & 18 & 0.002 & [-0.004, 0.007] & 0.524 & est. \\
Medical waiting & Unweighted CC FE & 282 & 30 & 10 & -0.010 & [-0.055, 0.035] & 0.671 & est. \\
Medical waiting & Pop.-weighted CC FE & 282 & 30 & 10 & -0.004 & [-0.079, 0.072] & 0.927 & est. \\
Medical waiting & Balanced FE & 238 & 25 & 10 & -0.001 & [-0.045, 0.043] & 0.973 & est. \\
Medical waiting & IPW CC FE & 282 & 30 & 10 & -0.009 & [-0.056, 0.038] & 0.704 & est. \\
Medical waiting & MAR MI full target & 608 & 37 & 18 & -0.001 & [-0.041, 0.039] & 0.966 & est. \\
\end{longtable}
\endgroup

\vspace{0.25em}

\begin{minipage}{0.96\textwidth}
\scriptsize \textit{Source:} Author's estimates from the thesis regression and imputation pipeline. Observed-data cells use country-and-year FE where feasible; MAR MI cells pool imputed outcome-observed panels using Rubin's rules \citep{rubin_1987,little_rubin_2019}.
\end{minipage}
\end{table}

\begin{figure}[H]
\centering
\includegraphics[width=0.94\textwidth]{outcome_estimand_coefficient_stability}
\caption{Coefficient stability across outcomes and estimands. Points show estimated poverty/social-exclusion coefficients for rows that meet the feasibility rule; horizontal intervals show 95\% confidence intervals. Infeasible estimands are omitted from the plot but retained in Table~\ref{tab:outcome_estimand_coefficient_matrix}.}
\label{fig:outcome_estimand_coefficient_stability}
\vspace{0.25em}

\begin{minipage}{0.94\textwidth}
\scriptsize \textit{Source:} Author's estimates from the outcome-by-estimand matrix. The figure is a stability diagnostic for aggregate coefficients, not evidence of individual-level effects.
\end{minipage}
\end{figure}

\begin{table}[H]
\centering
\caption{Drift-stability summary for the poverty/social-exclusion stress-test coefficient. The table reports framework metrics generated from the outcome-by-estimand matrix; TDI summarizes target movement and CSI summarizes conclusion agreement under the declared tolerance rule.}
\label{tab:drift_stability_summary}
\begingroup
\scriptsize
\setlength{\tabcolsep}{3pt}
\begin{tabularx}{\linewidth}{@{}p{0.15\linewidth}p{0.08\linewidth}rrrrrp{0.08\linewidth}p{0.16\linewidth}p{0.14\linewidth}@{}}
\toprule
Outcome & Denom. & Families & Median $\theta$ & $\theta$ IQR & Dir. agree & CSI & Median TDI & Label & Rule reason \\
\midrule
Dental population & population & 5 & 0.377 & 0.077 & 0.800 & 0.231 & 0.342 & missingness-dependent & CC--MI gap=0.365 \\
Medical cost & population & 5 & 1.133 & 0.466 & 1.000 & 0.000 & 0.343 & missingness-dependent & CC--MI gap=1.037 \\
Medical distance & population & 5 & 0.281 & 0.187 & 1.000 & 0.000 & 0.343 & missingness-dependent & CC--MI gap=0.225 \\
Medical population & population & 5 & 0.479 & 0.105 & 0.800 & 0.000 & 0.343 & missingness-dependent & CC--MI gap=0.512 \\
Medical waiting & population & 5 & -0.019 & 0.041 & 1.000 & 0.590 & 0.342 & target-dependent & CSI=0.590 and dir. agree=1.000 \\
Dental need & same\_needs & 1 & -0.309 & NA & NA & NA & 0.000 & non-identifiable & <2 estimand families \\
Medical need & same\_needs & 1 & 0.229 & NA & NA & NA & 0.000 & non-identifiable & <2 estimand families \\
\bottomrule
\end{tabularx}
\endgroup

\vspace{0.25em}

\begin{minipage}{0.96\textwidth}
\scriptsize \textit{Source:} Author's calculations from \texttt{target\_drift\_components.csv}, \texttt{drift\_stability\_theta.csv}, and \texttt{conclusion\_stability.csv}. Labels are deterministic diagnostics from the generated framework metrics; they are not causal or policy-effect classifications.
\end{minipage}
\end{table}

\section{Framework-Parameter Sensitivity}

The drift-stability labels depend on declared diagnostic choices, especially the CSI tolerance \(\delta\). Table~\ref{tab:framework_parameter_sensitivity_summary} varies \(\delta\) from 0.05 to 0.20 and recomputes the labels. It also recomputes TDI under equal, support-focused, and balance-focused component weights. These checks do not prove that the thresholds are optimal. They show whether the reported audit labels are fragile to declared diagnostic choices. Only medical waiting changes at the loosest \(\delta=0.20\); TDI weighting does not change labels because classification uses coefficient stability, complete-case--MI gaps, observed-estimand movement, and feasibility rules rather than a TDI cutoff.

\begin{table}[H]
\centering
\caption{Framework-parameter sensitivity of final labels. CSI columns vary the tolerance parameter \(\delta\); the TDI column reports whether equal, support-focused, or balance-focused TDI weights alter any target-dependent label.}
\label{tab:framework_parameter_sensitivity_summary}
\begingroup
\scriptsize
\setlength{\tabcolsep}{3pt}
\begin{tabularx}{\linewidth}{@{}p{0.12\linewidth}p{0.06\linewidth}p{0.14\linewidth}p{0.11\linewidth}p{0.11\linewidth}p{0.11\linewidth}p{0.11\linewidth}p{0.08\linewidth}p{0.08\linewidth}@{}}
\toprule
Outcome & Denom. & Baseline & $\delta=.05$ & $\delta=.10$ & $\delta=.15$ & $\delta=.20$ & CSI change? & TDI change? \\
\midrule
Dental pop. & pop. & miss.-dep. & miss.-dep. & miss.-dep. & miss.-dep. & miss.-dep. & No & No \\
Med. cost & pop. & miss.-dep. & miss.-dep. & miss.-dep. & miss.-dep. & miss.-dep. & No & No \\
Med. dist. & pop. & miss.-dep. & miss.-dep. & miss.-dep. & miss.-dep. & miss.-dep. & No & No \\
Med. pop. & pop. & miss.-dep. & miss.-dep. & miss.-dep. & miss.-dep. & miss.-dep. & No & No \\
Med. wait & pop. & target-dep. & target-dep. & target-dep. & target-dep. & stable & Yes & No \\
Dental need & need & non-ident. & non-ident. & non-ident. & non-ident. & non-ident. & No & No \\
Med. need & need & non-ident. & non-ident. & non-ident. & non-ident. & non-ident. & No & No \\
\bottomrule
\end{tabularx}
\endgroup

\vspace{0.25em}

\begin{minipage}{0.96\textwidth}
\scriptsize \textit{Source:} Author's calculations from the generated drift-stability matrix. Detailed CSI values and TDI-weighted summaries are stored in \texttt{outputs/framework\_parameter\_sensitivity.csv}.
\end{minipage}
\end{table}

\FloatBarrier

\section{Semi-Synthetic Missingness Simulation}

The simulation starts from the 282-row complete primary-outcome reference panel and estimates the poverty/social-exclusion reference estimand before imposing MCAR, early-year, country-group, financing-block, MAR-style, MNAR high-poverty/high-unmet, and Eurostat-realistic missingness. It then compares complete-case FE, population-weighted FE, IPW FE, PMM MI, and MNAR-shift MI. The exercise does not claim a definitive European coefficient; it tests how estimators behave when the complete artificial panel is known before deletion. Each draw also records the TDI components, coefficient error, sign reversal, interval coverage, wrong-conclusion rate, and failed replicates.

\begin{table}[H]
\centering
\scriptsize
\caption{Semi-synthetic missingness simulation summary. The reference estimand is the poverty/social-exclusion coefficient from the complete 282-row primary panel. Each mechanism uses 100 replicates; the table reports selected estimators to show bias, sign reversal, interval coverage of the reference estimand, and wrong-conclusion rates.}
\label{tab:simulation_main_summary}
\begin{tabular}{p{0.19\linewidth}p{0.18\linewidth}rrrrr}
\toprule
Mechanism & Estimator & Mean bias & Med. abs. error & Sign rev. & CI cov. & Wrong concl. \\
\midrule
Country group & CC FE & 0.0072 & 0.0271 & 0.0 & 1.0 & 0.0 \\
Country group & MNAR-shift MI & -0.1497 & 0.1476 & 0.15 & 0.01 & 0.15 \\
Country group & PMM MI & -0.1458 & 0.1466 & 0.12 & 0.02 & 0.12 \\
Covariate block & CC FE & -0.0048 & 0.0342 & 0.0 & 1.0 & 0.0 \\
Covariate block & MNAR-shift MI & -0.0028 & 0.0026 & 0.0 & 1.0 & 0.0 \\
Covariate block & PMM MI & -0.0026 & 0.0025 & 0.0 & 1.0 & 0.0 \\
Early-year & CC FE & -0.0773 & 0.0773 & 0.0 & 1.0 & 0.0 \\
Early-year & MNAR-shift MI & -0.2299 & 0.2303 & 1.0 & 0.0 & 1.0 \\
Early-year & PMM MI & -0.2304 & 0.2316 & 1.0 & 0.0 & 1.0 \\
Eurostat-realistic & CC FE & -0.0156 & 0.0439 & 0.0 & 0.94 & 0.0 \\
Eurostat-realistic & MNAR-shift MI & -0.0857 & 0.0879 & 0.0 & 0.46 & 0.0 \\
Eurostat-realistic & PMM MI & -0.0603 & 0.0612 & 0.0 & 0.79 & 0.0 \\
MAR observed & CC FE & -0.0019 & 0.0274 & 0.0 & 1.0 & 0.0 \\
MAR observed & MNAR-shift MI & -0.1471 & 0.1464 & 0.07 & 0.01 & 0.07 \\
MAR observed & PMM MI & -0.1419 & 0.1422 & 0.06 & 0.02 & 0.06 \\
MCAR & CC FE & 0.0028 & 0.1212 & 0.15 & 0.89 & 0.15 \\
MCAR & MNAR-shift MI & -0.1448 & 0.1434 & 0.05 & 0.0 & 0.05 \\
MCAR & PMM MI & -0.1315 & 0.1337 & 0.0 & 0.01 & 0.0 \\
MNAR high-pov./unmet & CC FE & -0.0164 & 0.0329 & 0.0 & 0.95 & 0.0 \\
MNAR high-pov./unmet & MNAR-shift MI & -0.1427 & 0.14 & 0.07 & 0.0 & 0.07 \\
MNAR high-pov./unmet & PMM MI & -0.1483 & 0.1484 & 0.12 & 0.01 & 0.12 \\
\bottomrule
\end{tabular}

\vspace{0.25em}

\begin{minipage}{0.96\textwidth}
\scriptsize \textit{Source:} Author's semi-synthetic simulation based on the complete-case Eurostat primary panel. The full validation output is generated in \texttt{outputs/simulation\_validation\_summary.csv}; the reference estimand is defined inside the simulation and is not a claim about a true European causal coefficient.
\end{minipage}
\end{table}

Table~\ref{tab:simulation_main_summary} shows that the direction of failure depends on the missingness mechanism. Complete-case FE is relatively stable under country-group, covariate-block, MAR-style, and Eurostat-realistic deletion in this reference panel, but MCAR deletion can create sign reversals because retained rows become small and noisy. PMM MI and MNAR-shift MI are not uniformly better. Figure~\ref{fig:simulation_tdi_validation_combined} links the simulation back to the framework: larger diagnostic drift is not a theorem of bias, but the generated draws show why target definition, row loss, estimator sensitivity, and non-identifiability should be reported before interpreting substantive coefficients.

\begin{figure}[H]
\centering
\includegraphics[width=0.98\textwidth]{simulation_tdi_validation_combined}
\caption{TDI validation diagnostics in the semi-synthetic simulation. Panel A compares mean TDI with median absolute coefficient error by mechanism--estimator cell. Panel B compares wrong-conclusion probability across TDI bins.}
\label{fig:simulation_tdi_validation_combined}
\vspace{0.25em}

\begin{minipage}{0.94\textwidth}
\scriptsize \textit{Source:} Author's semi-synthetic missingness simulation. TDI is used as an audit diagnostic; the figure does not claim that TDI proves bias.
\end{minipage}
\end{figure}

\begin{table}[htbp]
\centering
\small
\caption{Conceptual role of variable blocks in the aggregate panel. The table organizes interpretation for country-year associations and does not identify intervention effects.}
\label{tab:conceptual_variable_roles}
\begin{tabular}{@{}p{0.24\textwidth}p{0.30\textwidth}p{0.36\textwidth}@{}}
\toprule
Block & Variables & Role in interpretation \\
\midrule
Social vulnerability & Poverty or social exclusion, unemployment, inequality & Main associational exposure block; not individual-level risk. \\
Macroeconomic context & GDP per capita, GDP growth & Country-year resource context and confounding proxy. \\
Health-system financing & Government expenditure, compulsory financing, OOP share & Affordability and institutional proxies; not direct intervention estimates. \\
Capacity & Physicians and hospital beds & Availability proxies, mostly for prediction and sensitivity. \\
Country and year structure & Country indicators, year indicators & Persistent national differences and common period shocks. \\
Missingness selection & Inclusion probability, coverage flags, imputation predictors & Data-availability process that changes the analytical target. \\
\bottomrule
\end{tabular}
\vspace{0.25em}

\begin{minipage}{0.94\textwidth}
\scriptsize \textit{Source:} Author's synthesis from Chapter~\ref{ch:background_literature} and the variables in Table~\ref{tab:variable_summary}. The table defines interpretation roles for the aggregate model.
\end{minipage}
\end{table}

This conceptual table is not a causal identification graph. It clarifies why the same coefficient can change when the row set, adjustment block, or missing-data assumption changes; ecological aggregation still prevents individual-level interpretation \citep{robinson_1950,greenland_morgenstern_1989}.

\section{Complete-Case Selection}

The intended monitoring target contains 608 country-years with observed unmet medical care need. The selected fixed-effects regression sample contains 282 country-years complete for the baseline covariates. Table~\ref{tab:complete_case_selection_model} therefore comes before the baseline regression: it diagnoses how the analytical row set is formed.

The selection model is not causal. It asks whether inclusion in the complete-case regression is related to observed unmet need and observed macro conditions. Country-group terms are included in the diagnostic model but omitted from the compact displayed table. Predictable inclusion means the complete-case coefficient should be read as a selected-panel estimand, not as the coefficient for the full outcome-observed monitoring target.

\begin{table}[H]
\centering
\small
\caption{Complete-case inclusion model. The diagnostic uses all 608 outcome-observed country-years to model inclusion in the 282-row baseline complete-case regression sample; it is unweighted and country-clustered.}
\label{tab:complete_case_selection_model}
\begin{tabular}{@{}lrrr@{}}
\toprule
Predictor & Coef. & SE & p-value \\
\midrule
Unmet need & -0.000 & 0.007 & 0.951 \\
Log GDP & 0.145 & 0.070 & 0.038 \\
Unemployment & -0.016 & 0.006 & 0.009 \\
GDP missing & -0.524 & 0.135 & <0.001 \\
Unemployment missing & -0.279 & 0.049 & <0.001 \\
Year & 0.036 & 0.005 & <0.001 \\
\bottomrule
\end{tabular}

\vspace{0.2em}
\begin{minipage}{0.92\textwidth}\scriptsize Notes: Linear probability model for inclusion in the baseline complete-case regression sample among the 608 outcome-observed rows, with country-clustered standard errors. GDP and unemployment are median-imputed for this diagnostic and accompanied by missingness indicators.\end{minipage}

\vspace{0.25em}

\begin{minipage}{0.92\textwidth}
\scriptsize \textit{Source:} Author's calculations from the outcome-observed Eurostat primary panel. The model is a selection diagnostic following missing-data reporting logic \citep{little_rubin_2019,seaman_white_2013}.
\end{minipage}
\end{table}

\begin{figure}[H]
\centering
\includegraphics[width=0.82\textwidth]{missingness_inclusion_probability}
\caption{Estimated probability of complete-case inclusion among the 608 outcome-observed country-years. The figure compares rows included in the 282-row baseline complete-case FE model with outcome-observed rows excluded because covariates are missing.}
\label{fig:missingness_inclusion_probability}
\vspace{0.25em}

\begin{minipage}{0.9\textwidth}
\scriptsize \textit{Source:} Author's predictions from the complete-case inclusion model in Table~\ref{tab:complete_case_selection_model}.
\end{minipage}
\end{figure}

\FloatBarrier

\section{Selected Complete-Case Baseline}

The pooled specification relates unmet medical care need in country \(i\) and year \(t\) to country-year covariates and year indicators:
\[
y_{it} = \alpha + X_{it}\beta + \lambda_t + \varepsilon_{it}.
\]
The main fixed-effects specification adds country indicators:
\[
y_{it} = \alpha_i + X_{it}\beta + \lambda_t + \varepsilon_{it}.
\]
Both pooled and two-way fixed-effects models report country-clustered standard errors with a finite-sample correction, reflecting repeated observations within countries \citep{arellano_1987,bertrand_duflo_mullainathan_2004}.

\begin{table}[htbp]
\centering
\small
\caption{Baseline regression estimates for unmet medical care need. The FE column is the selected 282-row complete-case country-year sample over 2015--2024 with country and year fixed effects and country-clustered standard errors.}
\label{tab:regression_main_results}
\begin{tabular}{lcc}
\toprule
Variable & Pooled + year & Country + year \\
\midrule
Log GDP per capita & -0.546 (0.593) & -1.041 (2.201) \\
Unemployment rate & 0.160 (0.186) & -0.034 (0.087) \\
Poverty or social exclusion & 0.062 (0.065) & 0.177 (0.067) \\
Government health expenditure & 0.159 (0.190) & -0.229 (0.446) \\
Compulsory health financing & -0.401 (0.243) & 0.097 (0.444) \\
\midrule
Observations & 282 & 282 \\
Countries & 30 & 30 \\
Year indicators & Yes & Yes \\
Country indicators & No & Yes \\
$R^2$ & 0.221 & 0.895 \\
\bottomrule
\end{tabular}

\vspace{0.2em}
\begin{minipage}{0.92\textwidth}
\scriptsize Notes: Cells report coefficient estimates with standard errors in parentheses. Both the pooled model and the country-and-year fixed-effects model use country-clustered standard errors with finite-sample correction. No star thresholds are used because the thesis emphasizes coefficient stability across estimands.
\end{minipage}

\vspace{0.25em}

\begin{minipage}{0.92\textwidth}
\scriptsize \textit{Source:} Author's estimates from the 282-row complete-case Eurostat panel. The table is the selected-panel baseline used for later sensitivity comparisons.
\end{minipage}
\end{table}

The baseline country-and-year fixed-effects model gives a positive poverty/social-exclusion coefficient in the selected complete-case panel. This is the starting point, not the thesis conclusion. It uses 282 rows from 30 countries, while the outcome-observed monitoring target has 608 rows from 37 countries.

\section{Reduced-Covariate Model Ladder}

Before moving to IPW and imputation, Table~\ref{tab:model_ladder_results} asks whether the poverty/social-exclusion coefficient changes because of adjustment choices, country-universe restriction, population weighting, or balanced-outcome coverage. Models A--F are estimated with country and year fixed effects where feasible. The ladder starts with poverty only, then adds macroeconomic context, health financing, capacity, and broader theory-driven controls. The final reduced high-coverage model keeps the social-vulnerability and macro-labour structure while avoiding the larger coverage cost of system-expanded covariates.

\begin{table}[htbp]
\centering
\caption{Reduced-covariate model ladder for the poverty/social-exclusion coefficient. Rows report country-and-year FE estimates where feasible; infeasible rows are retained rather than forced.}
\label{tab:model_ladder_results}
\begingroup
\scriptsize
\setlength{\tabcolsep}{3pt}
\renewcommand{\arraystretch}{1.12}
\begin{tabularx}{\textwidth}{@{}p{0.26\textwidth}Yrrrp{0.20\textwidth}l@{}}
\toprule
Spec. & Target/weighting & Rows & Countries & Years & Coef. [95\% CI] & p-value \\
\midrule
A: Poverty only & Full available; unweighted & 375 & 36 & 11 & 0.157 [0.050, 0.264] & 0.004 \\
B: Poverty + GDP + unemployment & Full available; unweighted & 360 & 34 & 11 & 0.127 [0.009, 0.244] & 0.035 \\
C: B + health financing & Full available; unweighted & 282 & 30 & 10 & 0.178 [0.046, 0.309] & 0.008 \\
D: C + capacity & Full available; unweighted & 233 & 27 & 9 & 0.112 [-0.048, 0.272] & 0.171 \\
E: Full theory model & Full available; unweighted & 231 & 27 & 9 & 0.108 [-0.053, 0.268] & 0.189 \\
F: Reduced high-coverage model & Full available; unweighted & 358 & 34 & 11 & 0.132 [0.014, 0.249] & 0.028 \\
C-EU27: EU-27 baseline & EU-27; unweighted & 258 & 27 & 10 & 0.179 [0.034, 0.324] & 0.015 \\
C-PW: Population-weighted baseline & Full available; population-weighted & 282 & 30 & 10 & 0.304 [-0.007, 0.615] & 0.055 \\
C-BO: Balanced-outcome baseline & Balanced outcome; unweighted & 238 & 25 & 10 & 0.154 [0.060, 0.248] & 0.001 \\
C-BO-PW: Balanced-outcome weighted & Balanced outcome; population-weighted & 238 & 25 & 10 & 0.147 [-0.026, 0.320] & 0.096 \\
\bottomrule
\end{tabularx}

\vspace{0.2em}
\begin{minipage}{0.94\textwidth}\scriptsize Notes: All rows use country-and-year fixed effects and country-clustered standard errors. The ladder separates covariate adjustment from sample loss: A--F add or remove covariate blocks, C-EU27 changes country universe, C-PW changes weighting, and C-BO rows restrict to the balanced-outcome country set. Population weights are year-normalized country populations.\end{minipage}
\endgroup

\vspace{0.25em}

\begin{minipage}{0.96\textwidth}
\scriptsize \textit{Source:} Author's estimates from the Eurostat panel. Each row changes covariate set, weighting, country universe, or coverage restriction, so the table diagnoses coefficient stability across estimands.
\end{minipage}
\end{table}

\begin{figure}[htbp]
\centering
\includegraphics[width=0.92\textwidth]{model_ladder_coefficient_forest}
\caption{Coefficient stability across the model ladder. Points show the poverty/social-exclusion coefficient and 95\% confidence interval for feasible country-and-year FE specifications; rows differ by covariate set, weighting, country universe, or balanced-outcome restriction.}
\label{fig:model_ladder_coefficient_forest}
\vspace{0.25em}

\begin{minipage}{0.94\textwidth}
\scriptsize \textit{Source:} Author's estimates from Table~\ref{tab:model_ladder_results}. The plot visualizes how the same stress-test coefficient moves across model and target definitions.
\end{minipage}
\end{figure}

The model ladder separates three sources of movement. First, adding macroeconomic and financing variables changes the coefficient while also changing the retained country-year sample. Second, population weighting changes the monitoring question from equal-country to resident-weighted interpretation. Third, the EU-27 row tests whether the main selected-panel association is specific to the broader Eurostat-available country universe. These are estimand changes, not minor technical variants.

\section{Sensitivity Ladder}

Table~\ref{tab:main_sensitivity_ladder} is the main result table. It follows the poverty/social-exclusion coefficient across the model hierarchy, population weighting, balanced-outcome restriction, country-universe restrictions, lagged covariates, measurement-period checks, crisis-period exclusion, conservative country trends, IPW, MAR MI, and MNAR delta-shift sensitivity. Specifications with too few rows, countries, or years are reported as infeasible instead of forced.

The ladder shows the central empirical pattern. Observed covariate-complete specifications often remain positive, especially the country-and-year FE, balanced-outcome, EU-27, and IPW rows. The lagged-covariate, 2015-exclusion, and country-trend rows weaken precision. The MAR and MNAR imputation rows move the estimand to the full outcome-observed target and attenuate the poverty coefficient toward zero. The coefficient is therefore not a property of ``Europe'' in general; it is a property of a defined estimand.

The EU-27 row is close to the full Eurostat-available complete-case estimate, while smaller EEA and candidate-country universes are not large enough for stable fixed-effects inference. This supports the interpretation that the main difference is not EU-27 versus full operational Europe, but selected complete cases versus broader outcome-observed targets.

\begin{table}[H]
\centering
\scriptsize
\caption{Main sensitivity ladder for the poverty/social-exclusion coefficient. Rows use the generated thesis pipeline and distinguish model hierarchy, target population, universe, weighting, selection adjustment, and imputation assumptions.}
\label{tab:main_sensitivity_ladder}
\resizebox{\textwidth}{!}{%
\begin{tabular}{@{}p{0.18\textwidth}p{0.23\textwidth}rrp{0.16\textwidth}p{0.22\textwidth}@{}}
\toprule
Step & Target & Countries & Rows & Coef. [95\% CI] & Interpretation \\
\midrule
Pooled complete-case & selected complete-case rows & 30 & 282 & 0.062 [-0.065, 0.189] & associational baseline before country FE \\
Country FE & selected complete-case rows & 30 & 282 & 0.234 [0.118, 0.351] & absorbs persistent country differences \\
Country + year FE & selected complete-case rows & 30 & 282 & 0.178 [0.046, 0.309] & selected-panel primary FE estimate \\
Population-weighted FE & selected complete-case rows & 30 & 282 & 0.304 [-0.007, 0.615] & resident-weighted selected-panel estimand \\
Balanced-outcome FE & balanced outcome/population country set & 25 & 238 & 0.154 [0.060, 0.248] & coverage-restricted monitoring target \\
Full Eurostat-available FE & full Eurostat-available complete-case rows & 30 & 282 & 0.178 [0.046, 0.309] & full operational Europe universe \\
EU-27 FE & EU-27 complete-case rows & 27 & 258 & 0.179 [0.034, 0.324] & EU institutional universe sensitivity \\
Extended Europe FE & extended europe complete-case rows & 30 & 282 & 0.178 [0.046, 0.309] & additional country-universe sensitivity \\
EEA-country FE & eea-country complete-case rows & 2 & 14 & not estimated & additional country-universe sensitivity \\
Candidate FE & candidate complete-case rows & 0 & 0 & not estimated & additional country-universe sensitivity \\
Lagged covariates & outcome-observed rows with lag-complete covariates & 30 & 280 & 0.074 [-0.035, 0.182] & temporal-ordering check \\
Excluding 2015 & selected complete-case rows excluding 2015 & 30 & 252 & 0.099 [-0.018, 0.216] & measurement-break diagnostic \\
Excluding COVID years & selected complete-case rows excluding 2020--2022 & 30 & 195 & 0.230 [0.064, 0.396] & crisis-period sensitivity \\
Country trends & selected complete-case rows & 30 & 282 & 0.125 [-0.027, 0.278] & conservative over-control sensitivity \\
IPW & selected complete-case rows & 30 & 282 & 0.178 [0.046, 0.310] & selection diagnostic, not structural correction \\
MAR MI & 608 outcome-observed rows & 37 & 608 & 0.014 [-0.071, 0.099] & model-based broader target \\
MNAR pessimistic & 608 outcome-observed rows & 37 & 608 & 0.010 [-0.086, 0.107] & missing rows shifted toward worse macro-social conditions \\
MNAR optimistic & 608 outcome-observed rows & 37 & 608 & 0.010 [-0.086, 0.106] & missing rows shifted toward less severe macro-social conditions \\
\bottomrule
\end{tabular}%
}

\vspace{0.2em}
\begin{minipage}{0.96\textwidth}\scriptsize Notes: The ladder follows the poverty/social-exclusion coefficient as the model, country universe, weighting rule, selection adjustment, and imputation assumption change. Infeasible rows are retained when fixed-effects estimation has too few rows, countries, or years. MAR and MNAR rows are model-based sensitivity estimands, not corrected truth.\end{minipage}

\vspace{0.25em}

\begin{minipage}{0.96\textwidth}
\scriptsize \textit{Source:} Author's estimates from the thesis regression, IPW, MAR MI, and MNAR sensitivity pipeline. Observed-data FE rows use country-clustered uncertainty; imputed rows are pooled sensitivity estimands.
\end{minipage}
\end{table}

\begin{table}[H]
\centering
\small
\caption{Country-universe sensitivity for the poverty association. Rows use unweighted country-and-year fixed-effects models on available covariate-complete rows within each country universe; infeasible rows are reported rather than forced.}
\label{tab:universe_sensitivity_results}
\begingroup
\small
\begin{tabular}{@{}lrrll@{}}
\toprule
Universe & Countries & Rows & Coef. [95\% CI] & Status \\
\midrule
Full & 30 & 282 & 0.178 [0.046, 0.309] & estimated \\
EU-27 & 27 & 258 & 0.179 [0.034, 0.324] & estimated \\
EEA countries & 2 & 14 &  & infeasible (<50 rows) \\
Extended & 30 & 282 & 0.178 [0.046, 0.309] & estimated \\
Candidates & 0 & 0 &  & infeasible (<50 rows) \\
\bottomrule
\end{tabular}

\vspace{0.2em}
\begin{minipage}{0.92\textwidth}\scriptsize Notes: Rows report the poverty/social-exclusion coefficient from unweighted country-and-year fixed-effects models estimated on available covariate-complete rows within each country universe. Infeasible rows are not forced when the universe has too few complete-case countries, years, or rows.\end{minipage}
\endgroup

\vspace{0.25em}

\begin{minipage}{0.92\textwidth}
\scriptsize \textit{Source:} Author's estimates from country-universe filters defined in Chapter~\ref{ch:data_indicators}. Infeasible rows are retained to show support limits.
\end{minipage}
\end{table}

\FloatBarrier

\section{Inference Diagnostics}

The selected complete-case coefficient is also checked against alternative inference treatments. Table~\ref{tab:fe_inference_comparison} compares country-clustered, Driscoll--Kraay, and wild-cluster bootstrap inference for the same 282-row country-and-year FE estimand.

\begin{table}[htbp]
\centering
\small
\caption{Inference diagnostics for the selected complete-case poverty coefficient. The table compares country-clustered, Driscoll--Kraay, and wild-cluster bootstrap inference for the same 282-row country-and-year FE estimand.}
\label{tab:fe_inference_comparison}
\begin{tabular}{@{}lrrrrr@{}}
\toprule
Term & Coef. & Cluster SE & DK SE & Wild \(p\) & Wild 95\% CI \\
\midrule
Poverty/social exclusion & 0.178 & 0.067 & 0.076 & 0.004 & [0.058, 0.297] \\
\bottomrule
\end{tabular}

\vspace{0.2em}
\begin{minipage}{0.94\textwidth}\scriptsize Notes: The coefficient is from the selected complete-case country-and-year fixed-effects model. Cluster SE uses country clusters with finite-sample correction; DK is Driscoll--Kraay with one lag; the wild-cluster bootstrap uses country-level Rademacher weights and reports the bootstrap \(p\)-value and percentile interval.\end{minipage}

\vspace{0.25em}

\begin{minipage}{0.94\textwidth}
\scriptsize \textit{Source:} Author's inference diagnostics for the selected complete-case FE model. The comparison addresses dependence and small-cluster sensitivity but does not alter the target population.
\end{minipage}
\end{table}

\begin{table}[htbp]
\centering
\small
\caption{Fit diagnostics for the selected complete-case FE model. The table foregrounds within-\(R^2\), country count, year coverage, and years per country so the absorbed country/year structure is not mistaken for explanatory power of the time-varying covariates.}
\label{tab:fe_fit_diagnostics}
\resizebox{\textwidth}{!}{%
\begin{tabular}{@{}lrrrrrrrrr@{}}
\toprule
Model & \(N\) & Countries & Years & Avg. years & Min years & Max years & Within \(R^2\) & Between \(R^2\) & Overall \(R^2\) \\
\midrule
Baseline FE & 282 & 30 & 10 & 9.4 & 6 & 10 & 0.098 & 0.229 & 0.895 \\
\bottomrule
\end{tabular}%
}

\vspace{0.2em}
\begin{minipage}{0.94\textwidth}\scriptsize Notes: The baseline fixed-effects sample is the complete-case panel used for the selected poverty coefficient. The within \(R^2\) is emphasized because country and year effects absorb much of the overall variation.\end{minipage}

\vspace{0.25em}

\begin{minipage}{0.94\textwidth}
\scriptsize \textit{Source:} Author's diagnostics from the baseline complete-case FE fit.
\end{minipage}
\end{table}

These diagnostics strengthen the presentation of the selected-panel estimate without changing its interpretation. The high overall \(R^2\) mostly reflects country and year indicators; within-\(R^2\) is much smaller. Robust inference diagnostics address serial dependence and the small number of country clusters, but they do not solve target-population drift.

\section{IPW Diagnostic}

IPW is retained as a diagnostic, not as a fix for structural Eurostat non-coverage \citep{seaman_white_2013}. The weights are estimated from observed unmet need, year, country group, macro conditions, missingness indicators, and coverage flags, then stabilized and truncated. In this dataset IPW changes little because the complete-case sample is selected mainly through structural covariate availability. Reweighting observed complete cases cannot recover rows where required covariates are absent.

\begin{table}[H]
\centering
\small
\caption{Missingness-aware comparison of the poverty association. Complete-case and IPW rows use the selected FE sample with country-clustered SE; MI rows pool estimates over imputed outcome-observed panels.}
\label{tab:missingness_robustness_results}
\begin{tabular}{@{}lrrrr@{}}
\toprule
Specification & Coef. & SE & p-value & \(N\) \\
\midrule
Complete-case two-way FE & 0.1775 & 0.0672 & 0.008 & 282 \\
Population-weighted complete-case FE & 0.3041 & 0.1585 & 0.055 & 282 \\
Earlier MI sensitivity & 0.0001 & 0.0012 & 0.899 & 422 \\
IPW two-way FE & 0.1782 & 0.0674 & 0.008 & 282 \\
PMM MI full target & 0.0139 & 0.0434 & 0.748 & 608 \\
\bottomrule
\end{tabular}

\vspace{0.2em}
\begin{minipage}{0.92\textwidth}
\scriptsize Notes: All rows report the poverty or social exclusion coefficient from country-and-year fixed-effects models. Complete-case and IPW rows use country-clustered standard errors. IPW weights are stabilized and truncated at the 1st and 99th percentiles. The improved multiple-imputation row pools 30 imputed full outcome-observed panels using Rubin's rules.
\end{minipage}

\vspace{0.25em}

\begin{minipage}{0.92\textwidth}
\scriptsize \textit{Source:} Author's estimates from the missingness-robustness pipeline. The table contrasts selected complete cases, IPW diagnostics, and MI sensitivity targets.
\end{minipage}
\end{table}

Population-weighted estimates use country-year total population from Eurostat \texttt{demo\_pjan}, normalized within year. They answer a resident-weighted monitoring question, not a universally more correct version of the country-weighted question.

\clearpage

\begin{table}[H]
\centering
\scriptsize
\caption{Target-population sensitivity of the poverty association. Rows compare unweighted, year-normalized population-weighted, balanced-outcome, and MI full-target estimands; complete-case rows use country-clustered FE estimates.}
\label{tab:target_population_sensitivity_results}
\resizebox{\textwidth}{!}{%
\begin{tabular}{@{}lllcrrp{0.22\textwidth}lp{0.17\textwidth}@{}}
\toprule
Spec. & Universe & Target & Weight & Countries & Rows & Coef. [95\% CI] & p-value & Reading \\
\midrule
CC & Full available & Selected CC & Unweighted & 30 & 282 & 0.178 [0.046, 0.309] & 0.008 & selected panel \\
EU-27 CC & EU-27 & Selected CC & Unweighted & 27 & 258 & 0.179 [0.034, 0.324] & 0.015 & EU-27 selected panel \\
CC pop. & Full available & Selected CC & Population & 30 & 282 & 0.304 [-0.007, 0.615] & 0.055 & selected resident-weighted \\
Balanced & Full available & Balanced outcome & Unweighted & 25 & 238 & 0.154 [0.060, 0.248] & 0.001 & observed-panel restriction \\
Balanced pop. & Full available & Balanced outcome & Population & 25 & 238 & 0.147 [-0.026, 0.320] & 0.096 & weighted restriction \\
MI full & Full available & Full observed & Unweighted & 37 & 608 & 0.014 [-0.071, 0.099] & 0.748 & model-based sensitivity \\
\bottomrule
\end{tabular}%
}

\vspace{0.2em}
\begin{minipage}{0.96\textwidth}\scriptsize Notes: The coefficient is for poverty or social exclusion in country-and-year fixed-effects models with country-clustered standard errors, except the MI row, which pools fixed-effects estimates using Rubin's rules. Full available means all Eurostat-available national units in the outcome-observed panel; EU-27 restricts to current EU member states. Population weights are country-year total population from \texttt{demo\_pjan}, normalized within year.\end{minipage}

\vspace{0.25em}

\begin{minipage}{0.96\textwidth}
\scriptsize \textit{Source:} Author's estimates from the Eurostat panel and \texttt{demo\_pjan}. The rows change the target population or weighting rule and therefore change the estimand being interpreted.
\end{minipage}
\end{table}

\clearpage

\section{MAR Multiple Imputation}

The MAR MI analysis asks whether the selected complete-case poverty coefficient survives when the target is broadened to the 608 outcome-observed rows under a bounded imputation design \citep{rubin_1987,little_rubin_2019}. The implementation uses predictive mean matching: missing covariate values are filled by donor values selected from observed cases with similar model-predicted values. Raw GDP and log GDP are not imputed independently; GDP per capita is derived from imputed log GDP when the raw value is missing.

\begin{table}[htbp]
\centering
\scriptsize
\caption{Baseline PMM diagnostics for imputed covariates. Diagnostics refer to the 608-row outcome-observed target, the 30-imputation baseline PMM model, and bounded plausible ranges for imputed covariates.}
\label{tab:mi_diagnostics_summary}
\resizebox{\textwidth}{!}{%
\begin{tabular}{@{}lrrrrr@{}}
\toprule
Variable & Missing & Bounds & Viol. & Obs. mean & Imp. mean \\
\midrule
Log GDP & 0.7\% & [7.0, 12.3] & 0 & 10.07 & 9.67 \\
Unemployment & 8.6\% & [0.0, 40.0] & 0 & 8.65 & 7.53 \\
Poverty/social exclusion & 38.3\% & [0.0, 60.0] & 0 & 22.83 & 24.81 \\
Government health expenditure & 17.4\% & [0.0, 15.0] & 0 & 6.28 & 5.58 \\
Compulsory financing & 27.0\% & [0.0, 15.0] & 0 & 6.52 & 5.55 \\
Gini & 32.7\% & [0.0, 60.0] & 0 & 30.26 & 30.29 \\
Physicians & 29.1\% & [0.0, 1000.0] & 0 & 358.74 & 355.50 \\
Hospital beds & 14.1\% & [0.0, 1500.0] & 0 & 488.33 & 480.61 \\
OOP share & 27.0\% & [0.0, 60.0] & 0 & 20.65 & 25.87 \\
Long-term unemployment & 8.9\% & [0.0, 30.0] & 0 & 3.85 & 3.33 \\
\bottomrule
\end{tabular}%
}

\vspace{0.2em}
\begin{minipage}{0.96\textwidth}\scriptsize Notes: Diagnostics refer to the baseline 30-imputation model over the 608 outcome-observed country-years. Bounds are the plausible ranges imposed before fitting the fixed-effects model. Viol. counts imputed values outside those bounds after clipping and should be zero.\end{minipage}

\vspace{0.25em}

\begin{minipage}{0.96\textwidth}
\scriptsize \textit{Source:} Author's MI diagnostics from the 608-row outcome-observed Eurostat target. The table reports plausibility checks for the imputation design, not substantive effects.
\end{minipage}
\end{table}

\begin{table}[htbp]
\centering
\small
\caption{Multiple-imputation sensitivity variants. Rows compare PMM variants with outcome inclusion, no-outcome specification, country indicators, lagged predictors, and a 50-imputation check for the poverty coefficient in the outcome-observed target.}
\label{tab:mi_variant_summary}
\begin{tabular}{@{}lrrrr@{}}
\toprule
Variant & Coef. & SE & p-value & \(M\) \\
\midrule
Baseline PMM MI & 0.0139 & 0.0434 & 0.748 & 30 \\
No-outcome PMM MI & -0.0308 & 0.0492 & 0.532 & 30 \\
Country-FE PMM MI & 0.2516 & 0.0896 & 0.005 & 30 \\
Lagged-predictor PMM MI & -0.0472 & 0.0566 & 0.405 & 30 \\
50-imputation PMM check & 0.0122 & 0.0432 & 0.778 & 50 \\
\bottomrule
\end{tabular}

\vspace{0.25em}

\begin{minipage}{0.92\textwidth}
\scriptsize \textit{Source:} Author's MI sensitivity estimates pooled over imputed outcome-observed panels. Disagreement across rows is interpreted as imputation-model dependence.
\end{minipage}
\end{table}

\begin{table}[htbp]
\centering
\small
\caption{PMM donor diagnostics for the baseline imputation matrix. The table reports observed and missing values available to the donor-matching step; warnings identify matrix columns with fewer than 20 possible donors.}
\label{tab:mi_pmm_donor_diagnostics}
\begin{tabular}{@{}lrrrr@{}}
\toprule
Matrix column & Observed & Missing & Donor pool & Warning \\
\midrule
log\_gdp\_per\_capita & 604 & 4 & 20 & no \\
unemployment\_rate\_pc & 556 & 52 & 20 & no \\
poverty\_or\_social\_exclusion\_pc & 375 & 233 & 20 & no \\
government\_health\_expenditure\_gdp\_pc & 502 & 106 & 20 & no \\
compulsory\_health\_financing\_gdp\_pc & 444 & 164 & 20 & no \\
gini\_income & 409 & 199 & 20 & no \\
physicians\_per\_100k & 431 & 177 & 20 & no \\
hospital\_beds\_per\_100k & 522 & 86 & 20 & no \\
oop\_health\_expenditure\_share\_pc & 444 & 164 & 20 & no \\
long\_term\_unemployment\_rate\_pc & 554 & 54 & 20 & no \\
\bottomrule
\end{tabular}

\vspace{0.2em}
\begin{minipage}{0.92\textwidth}\scriptsize Notes: Donor diagnostics describe the baseline PMM matrix before imputation. The effective donor pool is the smaller of \(k=20\) and the number of observed values for a matrix column. Warnings indicate columns where fewer than 20 observed donor values were available.\end{minipage}

\vspace{0.25em}

\begin{minipage}{0.92\textwidth}
\scriptsize \textit{Source:} Author's diagnostics from the PMM imputation matrix. Donor counts document whether each variable has enough observed values to support donor matching.
\end{minipage}
\end{table}

\begin{figure}[htbp]
\centering
\includegraphics[width=0.86\textwidth]{mi_observed_vs_imputed_baseline}
\caption{Observed and imputed distributions in the baseline MI diagnostic. The sample is the 608-row outcome-observed target; imputed values are shown only where baseline covariates are missing.}
\label{fig:mi_observed_vs_imputed_baseline}
\vspace{0.25em}

\begin{minipage}{0.9\textwidth}
\scriptsize \textit{Source:} Author's baseline PMM imputation diagnostics. The figure checks distributional plausibility of imputed covariates.
\end{minipage}
\end{figure}

\begin{figure}[htbp]
\centering
\includegraphics[width=0.72\textwidth]{mi_coefficient_distribution}
\caption{Distribution of the poverty/social-exclusion coefficient across the 30 baseline PMM imputations. The pooled estimate attenuates toward zero, but the distribution is interpreted as model-based sensitivity rather than as proof of a target-invariant coefficient.}
\label{fig:mi_coefficient_distribution}
\vspace{0.25em}

\begin{minipage}{0.86\textwidth}
\scriptsize \textit{Source:} Author's baseline PMM MI estimates. The figure shows between-imputation coefficient variation before Rubin pooling.
\end{minipage}
\end{figure}

The PMM variants expose model dependence. Baseline, no-outcome, and lagged-predictor variants attenuate the poverty coefficient toward zero, while the country-indicator variant is more positive. That pattern is not hidden: it is evidence that the imputed estimand depends on the imputation structure.

\section{MNAR Multiple Imputation}

MAR assumptions cannot be verified in Eurostat covariate coverage. Table~\ref{tab:mnar_sensitivity_results} therefore applies MNAR delta shifts only to originally missing values after the baseline PMM imputation. Poverty, GDP, and unemployment shifts are not corrected estimates. They are stress tests for whether the MAR attenuation is mechanically dependent on one favorable imputation model.

The MI attenuation is interpreted as model-based instability under a broader target population. It may reflect target-population sensitivity, imputation-model dependence, MAR violations, or structural Eurostat non-coverage. The imputation model itself is part of the sensitivity design. The MI estimates are not treated as corrected estimates. They test whether the complete-case poverty coefficient is stable under broader model-based targets; it is not.

\begin{table}[H]
\centering
\scriptsize
\caption{MNAR delta-shift sensitivity for the poverty coefficient. Scenarios shift only originally missing values after PMM imputation and are stress tests for MAR departures, not corrected estimates.}
\label{tab:mnar_sensitivity_results}
\resizebox{\textwidth}{!}{%
\begin{tabular}{@{}lrrrl@{}}
\toprule
Scenario & Coef. [95\% CI] & p-value & \(M\) & Interpretation \\
\midrule
GDP -10\% & 0.011 [-0.085, 0.107] & 0.823 & 30 & MNAR stress test \\
GDP -20\% & 0.012 [-0.085, 0.108] & 0.811 & 30 & MNAR stress test \\
GDP -5\% & 0.011 [-0.086, 0.107] & 0.828 & 30 & MNAR stress test \\
Optimistic combo & 0.010 [-0.086, 0.106] & 0.837 & 30 & MNAR stress test \\
Pessimistic combo & 0.010 [-0.086, 0.107] & 0.831 & 30 & MNAR stress test \\
Poverty +10 pp & 0.007 [-0.089, 0.103] & 0.890 & 30 & MNAR stress test \\
Poverty +2 pp & 0.010 [-0.086, 0.107] & 0.834 & 30 & MNAR stress test \\
Poverty +5 pp & 0.010 [-0.086, 0.106] & 0.845 & 30 & MNAR stress test \\
Unemployment +2 pp & 0.011 [-0.085, 0.107] & 0.827 & 30 & MNAR stress test \\
Unemployment +5 pp & 0.011 [-0.084, 0.107] & 0.817 & 30 & MNAR stress test \\
\bottomrule
\end{tabular}%
}

\vspace{0.2em}
\begin{minipage}{0.96\textwidth}\scriptsize Notes: MNAR scenarios apply documented shifts only to originally missing values after the PMM MAR imputation. They are stress tests for model dependence and should not be read as corrected estimates.\end{minipage}

\vspace{0.25em}

\begin{minipage}{0.96\textwidth}
\scriptsize \textit{Source:} Author's MNAR delta-shift sensitivity analysis. Shifts are applied only to cells originally missing before imputation.
\end{minipage}
\end{table}

\begin{table}[H]
\centering
\small
\caption{MNAR tipping-point grid for the poverty/social-exclusion coefficient. Shifts are applied only to originally missing poverty/social-exclusion values after PMM imputation and are compared with the selected complete-case FE coefficient.}
\label{tab:mnar_tipping_point_results}
\begin{tabular}{@{}lrrrl@{}}
\toprule
Poverty shift & Coef. [95\% CI] & Distance & p-value & Status \\
\midrule
-10 pp & 0.019 [-0.077, 0.116] & 0.158 & 0.697 & grid point \\
-5 pp & 0.019 [-0.077, 0.116] & 0.158 & 0.697 & grid point \\
0 pp & 0.019 [-0.077, 0.116] & 0.158 & 0.697 & grid point \\
2 pp & 0.019 [-0.077, 0.115] & 0.158 & 0.699 & grid point \\
5 pp & 0.018 [-0.078, 0.114] & 0.159 & 0.712 & grid point \\
10 pp & 0.015 [-0.080, 0.110] & 0.162 & 0.756 & grid point \\
15 pp & 0.008 [-0.085, 0.102] & 0.169 & 0.865 & grid point \\
20 pp & -0.014 [-0.108, 0.081] & 0.191 & 0.779 & grid point \\
25 pp & -0.061 [-0.165, 0.042] & 0.239 & 0.246 & grid point \\
30 pp & -0.126 [-0.239, -0.013] & 0.303 & 0.029 & grid point \\
40 pp & -0.206 [-0.325, -0.087] & 0.384 & <0.001 & grid point \\
\bottomrule
\end{tabular}

\vspace{0.2em}
\begin{minipage}{0.92\textwidth}\scriptsize Notes: Shifts are applied only to originally missing poverty/social-exclusion cells after PMM imputation. The complete-case target coefficient is 0.178. Tipping-point status: not reached within grid; estimated shift: not reached within the grid.\end{minipage}

\vspace{0.25em}

\begin{minipage}{0.92\textwidth}
\scriptsize \textit{Source:} Author's MNAR tipping-point analysis. The table asks how far missing poverty/social-exclusion values would need to be shifted to approach the selected complete-case coefficient.
\end{minipage}
\end{table}

\begin{figure}[H]
\centering
\includegraphics[width=0.78\textwidth]{mnar_tipping_point_curve}
\caption{MNAR tipping-point curve for the poverty/social-exclusion coefficient. The horizontal dashed line is the selected complete-case FE coefficient; the curve shows pooled MI coefficients after shifting originally missing poverty/social-exclusion values.}
\label{fig:mnar_tipping_point_curve}
\vspace{0.25em}

\begin{minipage}{0.86\textwidth}
\scriptsize \textit{Source:} Author's MNAR tipping-point analysis. The curve visualizes imputation-model dependence under explicit departures from MAR.
\end{minipage}
\end{figure}

\FloatBarrier

\section{Interpretation}

The regression conclusion is estimand dependence, not a stable poverty effect. The selected complete-case model gives a substantively plausible positive association, and several observed covariate-complete perturbations point in the same direction. But the coefficient changes with model structure, weighting, measurement-period restrictions, country trends, and especially the move from selected complete cases to model-based outcome-observed targets.

This is the thesis spine. Complete-case modelling in open Eurostat health-monitoring panels can silently move the analyst from an intended European monitoring population to a selected covariate-complete panel. The poverty coefficient is useful because it reveals that drift; it is not presented as a durable causal, policy, or individual-level finding.
